% Preamble/Theorems/Theorems.tex

% Break style with line break and indentation after header
% Note: For examples that START with a list, always include introductory text first

\newtheoremstyle{break}
  {\topsep}{\topsep}%       space above and below
  {}%                       body font
  {}%                       indent amount (empty = no indent for header)
  {\bfseries}{}%            head font, punctuation after head
  {0pt}%                    space after head
  {\thmname{#1}\thmnumber{ #2}\thmnote{ (#3)}:\hfill\null\\\hspace*{\parindent}}%  head spec
\theoremstyle{break}

\theoremstyle{break}
\newtheorem{num}{Numerotare}[section]

\theoremstyle{break}
\newtheorem{definition}[num]{Definition}

\theoremstyle{break}
\newtheorem{proposition}[num]{Proposition}

\theoremstyle{break}
\newtheorem{theorem}[num]{Theorem}

\theoremstyle{break}
\newtheorem{corollary}[num]{Corollary}

\theoremstyle{break}
\newtheorem{lemma}[num]{Lemma}

\theoremstyle{break}
\newtheorem{example}[num]{Example}

\theoremstyle{break}
\newtheorem*{remark}{Remark}

\theoremstyle{break}
\newtheorem*{goal}{Goal}

\theoremstyle{break}
\newtheorem*{intuition}{Intuition}

% Compendium style for chapter summaries - key results without proofs
\newtheoremstyle{compendium}
  {\topsep}{\topsep}%
  {\small}%
  {}%
  {\bfseries\scshape}{}%
  {0pt}%
  {\thmname{#1}\thmnote{ --- #3}:\hfill\null\\\hspace*{\parindent}}%
\theoremstyle{compendium}
\newtheorem*{compendium}{Chapter Compendium}

\renewcommand{\qedsymbol}{QED}

% Redefine proof environment with break + indent formatting
\makeatletter
\renewenvironment{proof}[1][\proofname]
{%
    \par\pushQED{\qed}%
    \normalfont \topsep\p@\@plus\p@\relax
    \trivlist
    \item[\hskip\labelsep\bfseries\itshape#1:]\hfill\null\\\hspace*{\parindent}\ignorespaces
}
{%
    \popQED\endtrivlist\@endpefalse
}
\makeatother
